\section{Panduan Penulisan Pengantar}
\label{sec:Panduan_Penulisan_Pengantar}

Di dalam bagian pengantar ini perlu dijelaskan secara ringkas mengenai topik yang akan dikerjakan, rangkuman motivasi dari pemilihan topik ini, serta ringkasan alur dan isi dari dokumen C-251 ini. Pengantar yang baik harus ditulis dengan ringkas, padat dan jelas. Sangat direkomendasikan bagian pengantar ditulis seringkas mungkin tidak melebihi 1200 kata (2 halaman).

\textcolor{black}{Topik yang dipilih dalam \textit{Capstone Project} ini harus memenuhi \textbf{setidaknya aspek A8} di bawah ini beserta setidaknya minimal 2 aspek lainnya dari A1 - A7 sebagai berikut:}
\begin{enumerate}[label=A\arabic*.]
    \item melibatkan masalah teknis yang luas atau saling bertentangan,
    \item tidak memiliki solusi yang gamblang,
    \item solusi masalah tidak dicakup oleh standar dan kode program yang ada saat ini,
    \item melibatkan berbagai kelompok atau pemangku kepentingan,
    \item mencakup banyak bagian komponen atau sub-permasalahan,
    \item melibatkan berbagai disiplin ilmu, atau
    \item memiliki konsekuensi yang penting dalam berbagai konteks.
    \item \textbf{[Wajib]} Berdasarkan pada pengetahuan dan ketrampilan yang telah diperoleh di perkuliahan sebelumnya.
\end{enumerate}
Sebagai pengantar dokumen, mahasiswa harus mampu menjelaskan alasan kenapa topik yang dipilih memenuhi setidaknya 3 aspek tersebut di atas. Mahasiswa juga harus mampu menjelaskan bahwa dalam proyek capstone ini akan digunakan standar-standar keteknikan serta menghadapi kendala-kendala realistis, baik teknis maupun nonteknis, berdasarkan pengetahuan dan keterampilan yang telah diperoleh dari perkuliahan sebelumnya.

\section{Informasi Mengenai Ujian Dokumen C-251}
\label{sec:Ujian_Dokumen_C-251}
    
    Perlu diketahui bahwa dokumen C-251 akan diuji dalam sidang yang beranggotakan komite \textit{capstone} yang akan menentukan nilai akhir mata kuliah \textit{Capstone} 1. Dokumen C-251 merupakan penutup \textit{Capstone} 1 yang merupakan suatu karya tulis yang informatif yang ditujukan untuk meyakinkan pembaca dan komite \textit{capstone} bahwa apa yang akan dilakukan dalam \textit{Capstone Project} Saudara adalah layak atau berharga. Selain itu, \textit{Capstone} 1 berusaha memberikan jaminan bahwa apa yang akan dilakukan bisa diimplementasikan dalam waktu yang masuk akal (biasanya tiga sampai enam bulan).
    

\section{Panduan Penulisan Latar Belakang}
\label{sec:Panduan_Penulisan_Pengantar}

    \subsection{Latar Belakang}
    \label{subsec:Latar_Belakang}

    Tuliskan \textit{critical review} terhadap permasalahan-permasalahan yang masih bersifat umum pada bagian ini. Latar belakang harus ditulis secara mengerucut, yaitu dari permasalahan yang bersifat global menuju permasalahan yang lebih spesifik pada topik yang ingin Anda selesaikan atau cari solusinya, namun tetap belum mengarah ke aspek teknis. Anda perlu meyakinkan pembaca mengapa permasalahan ini signifikan atau penting untuk diselesaikan (pikirkan: \textit{apa dampaknya jika tidak diselesaikan?}). Dalam menuliskannya, Anda perlu merujuk pada referensi-referensi faktual seperti berita, artikel ilmiah, atau data relevan lainnya. Seluruh klaim harus didukung oleh referensi faktual berupa data statistik, laporan lembaga internasional atau nasional, berita terpercaya, atau artikel ilmiah yang valid. Perlu diingat bahwa bagian ini benar-benar fokus pada pembahasan permasalahan yang terjadi sehingga hindari membahas metode pada bagian ini.
    \vspace{6pt}

    \noindent
    Contoh (versi singkat):

    "Ketahanan pangan menjadi isu penting di tingkat global sejak tahun 2000, mengingat pertumbuhan jumlah penduduk yang pesat dan perubahan iklim yang memengaruhi ketersediaan sumber daya alam. Dalam konteks Indonesia, sektor perikanan, khususnya budidaya akuatik, berperan besar dalam menyediakan pangan yang bergizi dan meningkatkan perekonomian lokal hingga 38\%-45\%. Budidaya udang galah di perairan payau menjadi salah satu komoditas unggulan, namun menghadapi tantangan besar terkait fluktuasi kadar oksigen terlarut dan suhu air yang mempengaruhi kelangsungan hidup udang. Perubahan kondisi ini sering kali menyebabkan kematian massal pada udang, yang berdampak negatif pada hasil budidaya \cite{setiawan_ShrimpChallenge_2020}. Para petani tambak di wilayah pesisir sering kali kesulitan mengontrol dua variabel ini secara bersamaan, mengingat ketergantungan yang tinggi pada teknologi yang belum tersedia secara luas. Oleh karena itu, diperlukan solusi berbasis teknologi yang dapat memonitor dan mengatur kadar oksigen serta suhu secara efisien dalam lingkungan tambak, guna memastikan produktivitas budidaya yang optimal \cite{wahyuni_ChallengeAndImpact_2021, hidayat_FutureDirections_2019}."

    \subsection{Penjabaran Masalah Umum Menjadi Masalah Teknis} \label{subsec:Penjabaran_Masalah_Umum_Menjadi_Masalah_Teknis}

    Permasalahan umum yang sudah dijabarkan di atas kemudian dibahas lebih mendetail dengan mengarah ke aspek permasalahan teknis. Kemudian, mahasiswa harus mampu merumuskan masalah yang akan diselesaikan berupa pernyataan (problem statement), minimal tiga kendala di antara kendala berikut ini: aksesibilitas, estetika, kode atau etika, pembuatan, pembiayaan, ergonomi, ekstensibilitas, fungsionalitas, interoperabilitas, pertimbangan non-teknis seperti hukum, ekonomi, lingkungan hidup, keselamatan, sosial budaya, perundang-undangan, dan lain-lain, pemeliharaan, manufakturabilitas, pemasaran, kebijakan, jadwal, standar, keberlanjutan, atau kegunaan. Selanjutnya, mahasiswa juga harus mampu menjelaskan nama mata kuliah - mata kuliah atau praktikum yang mendukung proyek ini beserta alasannya. Untuk membantu mahasiswa menguraikan permasalahan secara benar, diberikan contoh sebagai berikut.
    \vspace{6pt}

    % \textit{
    “Permasalahan umum yang ditawarkan oleh komite capstone 2026 adalah ketahanan pangan dan budidaya akuatik. Pada proyek ini, tim memilih topik budidaya akuatik lebih spesifik pada budidaya udang galah di perairan payau. Menurut Bean Atkinson (2021) \cite{bean_shrimp_2021_DummyRef}, udang galah sangat rentan terhadap perubahan kadar oksigen terlarut di air dan perubahan suhu. Oleh karena itu, diperlukan suatu sistem yang mampu memonitor sekaligus mengatur kadar oksigen di tambak-tambak petani di pesisir.
    % }

    % \textit{
    Topik ini merupakan topik yang kompleks disebabkan oleh beberapa hal sebagai berikut:
    % }
    % \textit{
    \begin{enumerate}[label=\alph*.]
        \item Permasalahan oksigen terlarut dan perubahan suhu tidak mudah untuk diselesaikan karena melibatkan dua variabel yang berdiri sendiri dan dimungkinkan adanya hubungan secara matematis yang tidak linear. Hal ini menyebabkan pengendalian oksigen dan suhu secara bersama-sama memerlukan pemahaman yang mendalam tentang hubungan oksigen terlarut dan suhu. (Aspek A1)
        \item Dimungkinkan jumlah spesies, ukuran individu per spesies, jumlah makanan yang diberikan menyumbang tingkat oksigen terlarut yang perlu dipelajari lagi secara mendalam dengan menggunakan prinsip-prinsip sains yang sesuai. (Aspek A6)
        \item Untuk bisa memberikan solusi yang terbaik diperlukan kerjasama antara petani atau petambak tidak dengan pemerintah daerah yang tidak selalu berjalan dengan baik. (Aspek A4)
        \item Jumlah dana yang tersedia sangat terbatas. DTETI hanya menyediakan dana maksimum Rp. 500.000,- per mahasiswa sehingga diperlukan solusi yang berupa prototype. Pengujian solusi yang berupa prototype perlu menggunakan konsep-konsep ilmiah untuk menjamin ketercapaian solusi riilnya. (Aspek A7)
        \item Oksigen terlarut di air dan suhu di tambak petani perlu dimodelkan secara matematis untuk menghasilkan solusi terbaik. Hal ini memerlukan design ofexperiment yang tepat dan pemodelan dengan persamaan diferensial yang sesuai. Tim perlu menginvestigasi lebih lanjut berdasarkan kemampuan dan pemahaman anggota tim dalam membuat model matematika yang tepat. (Aspek A8)”
    \end{enumerate}
    % }

    \subsection{Kendala-Kendala yang Perlu Dipertimbangkan} 
    \label{subsec:Kendala-Kendala_yang_Perlu_Dipertimbangkan}

    Setelah mengidentifikasi permasalahan-permasalahan teknis, Anda perlu menguraikan berbagai faktor atau hambatan yang mungkin mempengaruhi penerapan solusi yang diusulkan. Tujuannya adalah agar pembaca memahami bahwa solusi tidak hanya dipilih berdasarkan efektivitas teknis saja, melainkan juga memperhitungkan aspek ekonomi, lingkungan, keselamatan, hukum, dan aspek-aspek relevan lainnya. Dalam menjelaskan kendala-kendala ini, mahasiswa sebaiknya:
    \begin{itemize}
        \item menyebutkan secara singkat tujuan solusi yang dirancang,
        \item menjelaskan mengapa mempertimbangkan kendala menjadi penting,
        \item mengelompokkan kendala-kendala dalam beberapa kategori utama (misalnya ekonomi, lingkungan, keselamatan, hukum/legalitas), dan
        \item memberikan penjelasan singkat untuk masing-masing kategori kendala.
    \end{itemize}
    \vspace{6pt}
    
    \noindent
    Contoh:
    
    "Untuk mendapatkan solusi yang tepat berkaitan dengan kematian ikan disebabkan oleh kurangnya oksigen terlarut dan perubahan suhu yang cepat di tambak petani Tim Capstone A\_23 “Bean” perlu memikirkan berbagai kendala yang menyertai implementasi dari solusi ini nantinya. Diantaranya:
    \begin{enumerate}[label=\alph*.]
        \item Ekonomi\\
        Desain solusi yang dibuat harus bisa meningkatkan perekonomian di daerah tambak. Jika solusi diimplementasikan jumlah kematian udang bisa turun sedikitnya 10\% sehingga meningkatkan penghasilan petani setidaknya 5\% dengan asumsi 5\% udang tidak terjual.
        \item Lingkungan\\ 
        Desain solusi harus ramah lingkungan, tidak boleh menghasilkan limbah cair maupun gas, dan tidak boleh mematikan vegetasi tambak dan lain-lain.
        \item Keselamatan\\
        Desain solusi harus teruji aspek keselamatannya. Jika menggunakan tenaga listrik maka harus dipastikan aman bagi manusia dan komunitas tambak dan lain-lain.
        \item Hukum dan Aspek Legal serta Perijinan\\
        Desain harus memenuhi aspek hukum dan legal. Desain solusi tidak boleh menggunakan frekuensi radio yang tidak legal. Selain itu desain harus mudah dari sisi perijinan dan lain-lain."
    \end{enumerate}

    \subsection{Mata Kuliah dan Kemampuan yang Mendukung} \label{subsec:Mata_Kuliah_dan_Kemampuan_yang_Mendukung}

    Uraikan minimal 10 mata kuliah atau praktikum pendukung yang akan dipakai di dalam \textit{Capstone Project} ini beserta alasannya. Berikut adalah contoh sederhana yang, namun, dalam implementasinya Anda perlu menguraikan alasannya dengan lebih jelas.
    \vspace{6pt}
    
    \noindent
    Contoh :

    "Mata kuliah dan kemampuan yang mendukung penyelesaian \textit{Capstone Project} ini antara lain sebagai berikut.
    \begin{enumerate}[label=\alph*.]
        \item Kalkulus Variabel Jamak\\
        Setiap udang akan berkontribusi terhadap kadar oksigen terlarut. Jika nantinya setiap udang dianggap sebagai satu variabel, maka untuk menghitung laju kenaikan oksigen terlarut diperlukan konsep kalkulus variabel jamak.
        \item Teori Vektor Matriks dan Aljabar Linear\\ 
        Teori vektor dan matriks akan digunakan untuk menganalisis hubungan linear antara jumlah udang, jumlah pakan dan kadar oksigen yang optimal menggunakan metode Singular Value Decomposition (SVD).
        \item Persamaan Diferensial\\
        Pertumbuhan udang secara dinamis perlu dimodelkan ke dalam persamaan diferensial. Tim akan membuat model pertumbuhan udang dari hari ke hari menggunakan pendekatan persamaan diferensial linear untuk melakukan prediksi sehingga bisa didapatkan solusi terbaik.
        \item Isyarat dan Sistem\\
        Jika desain menggunakan alat elektronika berupa sensor yang ditransmisikan dari tengah kolam maka memerlukan pemahaman tentang isyarat.
        \item Teknik Kendali Dasar dan Teknik Kendali Modern\\
        Salah satu solusi yang mungkin adalah pengendalian oksigen di kolam. Sehingga diperlukan teori teknik kendali.
        \item Pemrograman Berorientasi Objek\\ 
        Tim mahasiswa dari TIF perlu membuat sistem interface yang tapat dan mendesain algoritme yang kompleks oleh kerana itu diperlukan mata kuliah ini.
        \item Elektronika Analog\\
        Desain sistem instrumentasi dan sensor memerlukan dasar-dasar yang didapatkan dari mata kuliah elektronika analog.
        \item Pengukuran dan Instrumentasi\\
        Dalam membuat manual kerja prototipe, diperlukan konsep prinsip alat ukur listrik, menyebutkan dan menjelaskan alat ukur listrik analog dan perluasannya.
        \item Analisis Untai Elektrik DC\\
        Analisis rangkaian DC diperlukan untuk menghitung kebutuhan model elektronika dan nilai komponen R-C-L dan Op-Amp dengan menggunakan hukum-hukum dasar, misalnya Hukum Ohm dan Hukum Kirchhoff, analisis Nodal dan Mesh, Superposisi, Teorema Thevenin dan Norton.
        \item Probabilitas dan Variabel Acak\\
        Perlu menerapkan konsep variabel acak diskrit, menghitung distribusi probabilitas, fungsi densitas probabilitas, fungsi ekspektasi dan variansi dari suatu peubah dalam pengujian prototipe."
    \end{enumerate}

\subsection{Usulan Solusi Potensial} 
\label{subsec:Solusi Potensial}
Setelah memaparkan proyek/solusi-solusi terdahulu, penulis perlu mengidentifikasi kekurangan atau celah (gap) yang belum berhasil diselesaikan atau belum diimplementasikan di permasalahan yang diangkat. Paragraf ditutup dengan pernyataan jelas mengenai proyek capstone yang diusulkan, mencakup nama atau deskripsi singkat proyek, ukuran atau skala yang terukur seperti dimensi, kapasitas, maupun target spesifikasi, serta harapan atau target capaian yang ingin diraih dari proyek ini. Perlu diingat bahwa seluruh bagian latar belakang harus tetap fokus pada permasalahan, bukan pada metode atau solusi teknis. Penyebutan nama metode, algoritma, atau teknologi spesifik yang akan digunakan harus dihindari dan disimpan untuk bagian Metodologi.
